\section{Results}

All comparisons are paired $t$-tests across random seeds. Where a result is
described as \emph{replicated}, it was re-run on a seed set disjoint from
the one used to discover it.

\subsection{Graph structure improves classification, robustly}
Table~\ref{tab:periods} reports seven \emph{non-overlapping} evaluation
windows, each trained only on prior data and evaluated on the following
${\sim}15$ months. Graph structure improves per-date macro-F1 in
\textbf{7 of 7 windows} ($t=+11.67$), spanning the COVID drawdown, the 2022
bear market, and two recoveries. This is the most reproducible result in
the study and is not period-dependent.

\begin{table}[h]
\centering
\begin{tabular}{lccc}
\toprule
Evaluation window & macro-F1 $\Delta$ & tail-F1 $\Delta$ & LS spread $\Delta$ \\
\midrule
2016Q4--2017 & $+0.0209$ & $+0.0253$ & $-14.33$\,bp \\
2017Q4--2018 & $+0.0177$ & $+0.0192$ & $-12.09$\,bp \\
2018Q4--2019 & $+0.0218$ & $+0.0053$ & $-10.02$\,bp \\
2019Q4--2020 & $+0.0182$ & $+0.0019$ & $-12.46$\,bp \\
2020Q4--2021 & $+0.0163$ & $-0.0039$ & $-10.45$\,bp \\
2021Q4--2022 & $+0.0307$ & $+0.0048$ & $+5.69$\,bp \\
2022Q4--2024 & $+0.0216$ & $-0.0086$ & $-3.97$\,bp \\
\midrule
Graph wins & \textbf{7/7} & 5/7 & \textbf{1/7} \\
Mean ($t$, df$=6$) & $\mathbf{+0.0210}$ $(+11.67)$ & $+0.0063$ $(+1.38)$ & $\mathbf{-8.23}$\,bp $(-3.13)$ \\
\bottomrule
\end{tabular}
\caption{Sub-industry graph minus no-graph across seven independent windows,
3 seeds each. Classification improves everywhere; realised long-short spread
is worse almost everywhere.}
\label{tab:periods}
\end{table}

\subsection{Partition granularity: a step, not a gradient}
Table~\ref{tab:granularity} walks the graph partition down the GICS
hierarchy on \textbf{ten seeds disjoint from all other results}. Moving from
sector-level grouping to any finer partition improves tail-F1 by
${\approx}0.018$ ($p<.001$); beyond industry-group level the benefit
\emph{saturates}, with industry-group, industry and sub-industry
statistically indistinguishable. Notably this is the one effect in the
study that is \emph{larger and more significant on tail-F1 than on
macro-F1}, i.e.\ on the part of the distribution a portfolio actually uses.

\begin{table}[h]
\centering
\begin{tabular}{lcccc}
\toprule
Partition & Edges & macro-F1 & tail-F1 & vs.\ sector (tail-F1) \\
\midrule
Sector (11) & 25{,}782 & $0.3664$ & $0.2344$ & --- \\
Industry group (25) & 12{,}268 & $0.3728$ & $0.2524$ & $+0.0180$, $p<.01$ \\
Industry (67) & 5{,}904 & $0.3722$ & $0.2528$ & $+0.0184$, $p<.001$ \\
Sub-industry (124) & 3{,}548 & $0.3740$ & $0.2526$ & $+0.0183$, $p<.001$ \\
\bottomrule
\end{tabular}
\caption{Granularity, $n=10$ disjoint seeds. Finer partitions use
\emph{fewer} edges, so the mechanism is grouping precision, not connectivity.
The effect saturates by ${\sim}25$ groups.}
\label{tab:granularity}
\end{table}

An earlier version of this study reported a \emph{monotone} gradient with
fineness ($+0.0090$ macro-F1, $p<.05$, $n=5$). That did not replicate:
run-to-run nondeterminism alone moves the comparison between $p<.05$ and
non-significance. We report the step-function result because it survives
independent replication.

\subsection{Proprietary relationship data carries no premium}
With a genuine sector graph as the comparison, a graph built from Bloomberg
supplier/customer and institutional-holder relationships is statistically
indistinguishable from one built from free GICS labels
($\Delta = +0.0033$, $t=0.82$, n.s.), while both beat no-graph
($p<.001$ and $p<.01$). Supplier/customer edges recover ${\sim}90\%$ of the
proprietary graph's gain; institutional-holder overlap and rolling
return-correlation graphs add nothing ($\Delta = -0.006$ and $-0.005$, both
n.s.).

We note that an earlier version of this study reached the opposite
conclusion because its ``sector proxy'' baseline was, through a labelling
error, a random block graph. Against random edges the proprietary premium
appears significant ($p<.05$); against a real sector partition it vanishes.
We report this because validating relationship data against degenerate
proxies is a plausible failure mode elsewhere in this literature.

\subsection{The classification advantage does not become portfolio performance}
A Markowitz mean-variance layer converts predictions to portfolios
(long top-quintile / short bottom-quintile, dollar-neutral, gross cap 2.0,
5\,bps costs). On an adequately powered window (383 rebalances) the graph
signals underperform an identical no-graph portfolio: sub-industry Sharpe
$0.174$ vs.\ $1.117$ ($\Delta=-0.94$), Bloomberg $0.716$ ($\Delta=-0.40$).
The shortfall is not a transaction-cost effect --- the per-period
gross-return gap is $-44.4$\,bp against a cost gap of $+5.0$\,bp, so costs
explain ${\sim}10\%$. Neither signal shrinkage ($\lambda: 0 \to 0.95$,
Sharpe $0.431 \to 0.444$, n.s.) nor a transaction-cost-aware turnover
penalty (turnover $1.489 \to 1.340$ but Sharpe $0.438 \to 0.357$, $p<.05$)
recovers performance.

\paragraph{The sign is construction-dependent.} We emphasise that this
result is specific to mean-variance sizing. Re-running with equal-weight
quintile sizing and non-overlapping rebalances at the \emph{same} 5-day
horizon reverses it: the graph signal is marginally \emph{better}
($\Delta = +0.258$ Sharpe, $t=+1.99$). Consistent with
\citet{michaud1989markowitz}, equal-weight sizing of identical picks
exhibits $3.4\times$ lower seed variance than mean-variance sizing
($0.31$ vs.\ $1.07$), so at this effect size mean-variance comparisons
cannot establish direction. Every configuration --- including no-graph ---
concentrates ${\sim}50\%$ of gross exposure in a single sub-industry and
holds ${\sim}2$ effective groups out of 125, a property of gross-capped
mean-variance optimisation itself rather than of any signal.

We therefore do not claim the graph signal is economically harmful. What
the evidence supports is narrower and, we think, more interesting: a
robust, replicated classification gain yields \emph{no demonstrable}
portfolio advantage under any construction we tested.

\subsection{Horizon and holding period}
Table~\ref{tab:horizon} reports non-overlapping backtests at 5, 20 and 60
days (hold $H$ days, then trade), equal-weight quintiles, with Newey--West
standard errors on the paired difference.

\begin{table}[h]
\centering
\begin{tabular}{lccccc}
\toprule
Horizon & Signal & Sharpe & Mean/period & Turnover/yr & $n$ periods \\
\midrule
5d & graph & $0.368$ & $2.5$\,bp & $32.7$ & 101 \\
5d & no-graph & $0.110$ & $-2.3$\,bp & $20.4$ & 101 \\
20d & graph & $0.653$ & $19.3$\,bp & $11.9$ & 25 \\
20d & no-graph & $0.713$ & $34.5$\,bp & $7.0$ & 25 \\
60d & graph & $0.878$ & $91.2$\,bp & $4.7$ & 8 \\
60d & no-graph & $\mathbf{1.491}$ & $148.0$\,bp & $3.0$ & 8 \\
\midrule
\multicolumn{2}{l}{$\Delta$ Sharpe, 5d} & \multicolumn{4}{l}{$+0.258$ $(t{=}{+}1.99)$} \\
\multicolumn{2}{l}{$\Delta$ Sharpe, 20d} & \multicolumn{4}{l}{$-0.060$ $(t{=}{-}0.32)$} \\
\multicolumn{2}{l}{$\Delta$ Sharpe, 60d} & \multicolumn{4}{l}{$-0.614$ $(t{=}{-}1.22)$} \\
\bottomrule
\end{tabular}
\caption{Non-overlapping horizon backtests. Absolute performance rises and
annualised turnover falls ${\sim}7\times$ for both signals; the graph
advantage does not improve with horizon.}
\label{tab:horizon}
\end{table}

Longer holding periods improve absolute performance substantially for
\emph{both} signals while cutting annualised turnover roughly sevenfold ---
the cost channel identified above, now acting favourably. But the graph
advantage does not improve with horizon; it shrinks and changes sign.

Statistical power is the binding constraint here: non-overlapping 60-day
holds leave only eight periods per window--seed, so a Sharpe standard error
of roughly $1/\sqrt{8} \approx 0.35$ swamps the $0.614$ gap. Neither the
20-day nor the 60-day difference is significant, and we do not read the
point estimates as directional evidence. This is an intrinsic tension in
horizon selection: longer horizons raise Sharpe but destroy the sample
needed to verify it.

\paragraph{Industry rotation does not help.} Because the Fama--MacBeth
results below localise the signal between industries, we tested a
group-level rotation book (long top-quintile sub-industries, short bottom,
equal-weighting members within group) against the stock-level book on
identical predictions. It is significantly \emph{worse} (mean $-5.35$\,bp,
$t=-2.78$, $p<.05$; better in 1 of 7 windows). The inference from
between-industry identification to a group-level portfolio is too strong:
industry fixed effects shrinking a coefficient shows that between-industry
variation carries signal, not that within-industry variation carries none.
Aggregating to ${\sim}124$ group means discards real within-group
information and reduces effective bets from ${\sim}200$ stocks to
${\sim}50$ groups.

\subsection{Mechanism: the advantage lives where portfolios cannot reach}
Per-class F1 resolves the paradox (Table~\ref{tab:perclass}). Both graph
models' macro-F1 advantage comes almost entirely from the \textbf{Neutral}
class --- 60\% of the probability mass, and the one region a quintile
long-short book never trades --- while tail accuracy is flat-to-worse. The
realised long-short spread orders exactly as portfolio Sharpe does:
$6.0 < 7.4 < 9.7$\,bp against Sharpe $0.174 < 0.716 < 1.117$.

\begin{table}[h]
\centering
\begin{tabular}{lcccc}
\toprule
Signal & Down F1 & Neutral F1 & Up F1 & LS spread \\
\midrule
Sub-industry & $0.2579$ & $\mathbf{0.5852}$ & $0.2355$ & $6.0$\,bp \\
Bloomberg & $0.2118$ & $\mathbf{0.5795}$ & $0.2817$ & $7.4$\,bp \\
No graph & $0.2298$ & $0.5472$ & $0.2774$ & $\mathbf{9.7}$\,bp \\
\bottomrule
\end{tabular}
\caption{Per-class F1, long out-of-sample. Graph gains concentrate in
Neutral ($+0.032$ to $+0.038$); tail F1 is slightly worse ($-0.007$).}
\label{tab:perclass}
\end{table}

We tested and rejected three alternative mechanisms: transaction costs
(above); \emph{signal smoothing}, where within-group variance shares were
non-monotonic across configurations and inside noise; and \emph{portfolio
concentration}, where group Herfindahl indices were near-identical
($0.499$/$0.497$/$0.542$), with the no-graph book marginally the most
concentrated.

\subsection{Robustness}
The static supply-chain snapshot induces look-ahead whose severity scales
with distance from the pull date. If contamination drove the results, the
graph advantage should shrink as evaluation approaches the snapshot; it does
the opposite, being smallest in the most-contaminated window (2017,
$+0.0096$) and largest and tightest in the cleanest (2024, $+0.0192$,
$t=+19.7$). We therefore attribute the signal to persistent relationships,
for which a stale snapshot is a serviceable proxy.

Two quantities in this study are large relative to the effects measured:
run-to-run nondeterminism (which alone flips the granularity gradient
between $p<.05$ and n.s.) and evaluation-window choice (single-window
long-short spread estimates swing ${\sim}16$\,bp between runs of an
identical configuration). We accordingly report multi-window aggregates and
disjoint-seed replications throughout, and treat single-window or
single-seed comparisons as uninformative.

\subsection{The signal does not survive standard controls}
Table~\ref{tab:fm} reports Fama--MacBeth cross-sectional regressions of
forward 5-day returns on the standardised graph signal, with 12--2 momentum,
one-month reversal, 60-day volatility and GICS-sector fixed effects, averaged
across dates with Newey--West standard errors (lag 5, matching the overlapping
horizon). Even uncontrolled the return-predictive content is weak (pooled
$t=+1.74$, significant in three of seven windows and negative in two), so the
robust 7/7 macro-F1 result does not correspond to robust return prediction.
After controls the signal retains 23\% of its univariate magnitude and is not
significant.

\begin{table}[h]
\centering
\begin{tabular}{lcc}
\toprule
Evaluation window & Signal alone & $+$ controls $+$ industry FE \\
\midrule
2016Q4--2017 & $4.81$\,bp $(t{=}{+}2.57)$ & $2.75$\,bp $(t{=}{+}1.54)$ \\
2017Q4--2018 & $2.45$\,bp $(t{=}{+}0.88)$ & $-0.09$\,bp $(t{=}{-}0.04)$ \\
2018Q4--2019 & $3.48$\,bp $(t{=}{+}0.95)$ & $2.08$\,bp $(t{=}{+}0.93)$ \\
2019Q4--2020 & $9.10$\,bp $(t{=}{+}2.54)$ & $2.06$\,bp $(t{=}{+}0.83)$ \\
2020Q4--2021 & $-1.25$\,bp $(t{=}{-}0.16)$ & $-1.81$\,bp $(t{=}{-}0.33)$ \\
2021Q4--2022 & $-5.78$\,bp $(t{=}{-}0.87)$ & $-3.50$\,bp $(t{=}{-}0.95)$ \\
2022Q4--2024 & $5.30$\,bp $(t{=}{+}2.78)$ & $2.05$\,bp $(t{=}{+}2.15)$ \\
\midrule
\textbf{Pooled} & $\mathbf{2.88}$\,bp $(t{=}{+}1.74)$ & $\mathbf{0.67}$\,bp $(t{=}{+}0.64)$ \\
\bottomrule
\end{tabular}
\caption{Fama--MacBeth coefficients on the standardised graph signal,
bp per 5-day period, 707 cross-sectional regressions. Industry fixed effects
implemented as within-sector demeaning, so the coefficient is identified off
within-industry variation.}
\label{tab:fm}
\end{table}

Because the industry fixed effect is within-sector demeaning, this also
localises the signal: its predictive content is largely \emph{between}
industries, not within them. Message passing ranks whole peer groups well ---
which registers as macro-F1, concentrated in the Neutral class --- but does not
separate stocks inside a group, which is what a quintile book requires; and
group-level prediction is in turn largely subsumed by industry momentum.

Size and book-to-market are not controlled for, as neither is available in our
data. We note the direction of that limitation: the signal already fails
against a partial control set, so additional controls can only reduce it
further, not restore it.
